\section{Introduction}

The production of 2D animation, especially anime-style content, remains highly labor-intensive, requiring extensive manual refinement of motion across consecutive frames. As animation pipelines continue to expand across film, television, games, and digital media, automating motion-aware processing has become increasingly important. Optical flow is a fundamental component in this context, as it estimates dense pixel-wise correspondences between frames and supports a wide range of downstream tasks, including frame interpolation \cite{animeinterp,superSloMo,softmaxSplat}, line-art inbetweening \cite{linedrawingFlow}, temporally consistent color propagation \cite{toonTracking}, and coherence preservation in stylized video generation. Recent animation-oriented systems such as AnimeInterp \cite{animeinterp} further highlight that high-quality optical flow is essential for preserving the unique temporal structure of 2D animated content.

Although optical flow has been studied extensively in computer vision, most successful methods have been developed and evaluated for natural-image domains. Architectures such as PWC-Net \cite{pwcnet}, RAFT \cite{raft}, GMA \cite{gma}, and GMFlow \cite{gmflow} achieve strong performance on benchmarks such as Sintel \cite{sintel} and FlyingThings3D \cite{ft3d}. However, these methods often generalize poorly to 2D animation, whose visual statistics differ substantially from natural scenes. Anime frames typically contain large textureless regions, sharp contour boundaries, repeated colors, exaggerated non-rigid motion, and layered compositions with abrupt occlusion changes. As shown by the AnimeRun benchmark \cite{animerun}, optical flow models trained for natural imagery frequently drift in flat regions, blur motion boundaries, and exhibit unstable temporal behavior when applied to animation sequences.

A major challenge in building robust animation-specific optical flow systems is the scarcity of high-quality ground-truth motion annotations. Producing dense flow labels for stylized content requires controlled rendering pipelines, accurate handling of composited layers and occlusions, and costly manual verification. These requirements make supervised learning difficult to scale across diverse animation styles and long video sequences. For this reason, \emph{unsupervised optical flow learning} is particularly appealing for animation: it avoids reliance on expensive annotations and makes it possible to learn directly from raw animated videos.

However, existing unsupervised optical flow methods remain poorly suited to the animation domain. First, many methods rely heavily on photometric reconstruction losses, which become unreliable in flat-color cartoon regions with weak or repetitive texture. Second, most unsupervised models are designed for two-frame estimation, limiting their ability to exploit temporal context for stable motion reasoning. Third, they typically lack explicit structural priors, making it difficult to preserve object boundaries and region-level motion coherence in stylized scenes. Finally, their motion reasoning is often local or short-range, which reduces robustness under large character motion, fast camera movement, and long-range temporal dependencies.

To address these limitations, we propose \textbf{AniUnFlow-SamT}, an \emph{unsupervised segment-guided multi-frame optical flow framework for 2D animation with hybrid temporal transformers}. Our method is designed specifically for the challenges of animation videos. Instead of estimating motion from isolated image pairs, AniUnFlow-SamT processes short clips and performs multi-frame motion reasoning through global matching, latent temporal memory, and temporal regression. To better handle textureless regions and sharp motion discontinuities, we incorporate segment-guided structural priors derived from precomputed masks, which modulate correspondence estimation, inject segment-aware attention, and encourage object-level motion consistency. In addition, a hybrid refinement module iteratively sharpens flow predictions using boundary-aware guidance and uncertainty-aware optimization. The entire framework is trained without ground-truth optical flow using a combination of photometric reconstruction, smoothness regularization, forward-backward consistency, segment-level coherence, boundary-aware objectives, and multi-frame cycle constraints.

By combining unsupervised learning, multi-frame temporal reasoning, and segment-guided structure modeling, AniUnFlow-SamT is designed to produce more stable, coherent, and boundary-preserving optical flow for anime-style videos than conventional unsupervised two-frame approaches.

Recent optical flow advances also motivate this design direction. Large-scale correspondence modeling and iterative refinement have proven highly effective in modern flow estimators, yet most pipelines remain pixel-centric and do not explicitly enforce object-level motion coherence \cite{pwcnet,raft,flowformer}. In parallel, layered-motion studies suggest that piecewise transformations with visibility ordering are beneficial near occlusions and independently moving regions \cite{wang1994representing,sevillalara2016optical}. Foundation segmentation models such as SAM and SAM~2 further provide high-quality object partitions and video-aware segmentation priors \cite{sam,ravi2024sam2}. Together, these observations suggest that dense correspondence and object-level structure should be modeled as complementary priors.

\textbf{Our main contributions are as follows:}
\begin{itemize}
    \item We propose \textbf{AniUnFlow-SamT}, an unsupervised segment-guided multi-frame optical flow framework tailored to the unique motion and appearance characteristics of 2D animation.
    \item We introduce a hybrid temporal transformer architecture that combines global matching, latent temporal cost memory, and temporal regression to improve motion estimation beyond conventional two-frame unsupervised models.
    \item We incorporate segment-guided structural priors into unsupervised optical flow learning, enabling sharper motion boundaries, stronger intra-object consistency, and improved robustness in textureless animation regions.
    \item We design a fully unsupervised training objective that integrates photometric, smoothness, consistency, boundary-aware, segment-level, and multi-frame cycle constraints for learning animation-specific motion representations without ground-truth flow.
    \item We demonstrate the effectiveness of the proposed approach on anime-style optical flow benchmarks and real animation sequences, showing improved temporal stability and boundary quality in challenging large-motion scenarios.
\end{itemize}
